\documentclass[11pt]{article}

% --- PACKAGES ---
\usepackage[letterpaper, top=1in, left=1in, right=1in, bottom=1in]{geometry}
\usepackage{parskip}      % For paragraphs separated by vertical space instead of indentation
\usepackage{amsmath}      % For advanced math typesetting
\usepackage{fancyhdr}     % For custom headers and footers
\usepackage{tabularx}     % For tables with adjustable-width columns
\usepackage{longtable}    % For tables that span multiple pages
\usepackage{graphicx}     % For including images
\usepackage{xurl}         % For better URL line breaking
\usepackage[T1]{fontenc} % For font encoding
\usepackage{hyperref}     % For hyperlinks
\usepackage{booktabs}     % For better looking tables

% --- HYPERLINK SETUP ---
\hypersetup{
  colorlinks=true,
  linkcolor=blue,
  filecolor=magenta,
  urlcolor=cyan,
  pdftitle={STAT 579: Applied Causal Inference},
}

% --- COURSE INFORMATION COMMANDS ---
\newcommand{\coursenumber}{STAT 579}
\newcommand{\coursetitle}{Applied Causal Inference}
\newcommand{\semester}{Fall 2025}
\newcommand{\instructor}{Davood Tofighi, Ph.D.}
\newcommand{\email}{[Your UNM Email Address]}
\newcommand{\officehours}{[Day and Time], [Location or Zoom Link]}
\newcommand{\classmeetings}{Tuesdays and Thursdays, 9:30 AM – 10:45 AM}
\newcommand{\location}{[Building and Room Number]}

% --- HEADER AND FOOTER ---
\pagestyle{fancy}
\fancyhf{} % Clear all header and footer fields
\rhead{\semester}
\lhead{\coursenumber: \coursetitle}
\cfoot{\thepage}

\begin{document}

% --- TITLE ---
\begin{center}
  {\Huge \coursenumber: \coursetitle} \\
  \vspace{4mm}
  {\large Syllabus -- \semester}
\end{center}

\vspace{8mm}

% --- INFO TABLE ---
\noindent
\begin{tabularx}{\textwidth}{@{}ll}
  \textbf{Instructor:}     & \instructor \\
  \textbf{Affiliation:}    & Department of Mathematics and Statistics \\
  \textbf{Email:}          & \href{mailto:\email}{\email} \\
  \textbf{Office Hours:}   & \officehours \\
  \textbf{Class Meetings:} & \classmeetings \\
  \textbf{Class Location:} & \location \\
\end{tabularx}

\vspace{6mm}
\section*{Course Overview}
This graduate-level course provides a comprehensive introduction to the theory and application of causal inference, with a strong emphasis on biostatistical applications. We will move beyond association to learn formal frameworks for defining and estimating causal effects from observational data. The course covers foundational concepts like potential outcomes and causal graphs (DAGs), explores challenges such as confounding and selection bias, and details essential methods including propensity scores, instrumental variables, mediation analysis, and techniques for time-varying treatments.

\vspace{6mm}
\subsection*{My Approach to Causal Inference}
This course emphasizes that causal inference is a way of thinking rigorously about cause and effect. We will constantly return to foundational principles: clearly defining the causal question, making all assumptions explicit (often with DAGs), understanding what each method can and cannot do, and critically assessing the robustness of our conclusions.

\vspace{6mm}
\section*{Learning Objectives}
Upon successful completion of this course, students will be able to:
\begin{enumerate}
  \item Distinguish between association and causation and formally define causal effects using the potential outcomes framework.
  \item Represent causal structures using Directed Acyclic Graphs (DAGs) and use them to identify confounding and selection bias.
  \item Articulate and critically evaluate the assumptions required for causal identification (e.g., exchangeability, positivity, SUTVA, exclusion restriction).
  \item Apply propensity score methods (matching, weighting) to estimate causal effects, including checking diagnostics like balance and overlap.
  \item Implement g-computation and inverse probability weighting for estimating causal effects of time-varying treatments using Marginal Structural Models.
  \item Decompose effects into natural direct and indirect effects using causal mediation analysis, understanding the strong assumptions required.
  \item Conduct sensitivity analyses to evaluate the robustness of causal estimates to violations of key assumptions.
\end{enumerate}

\vspace{6mm}
\section*{Assignments and Grading Policies}

\vspace{6mm}
\subsection*{Assessment Components}
Your final grade is a reflection of your understanding of theoretical concepts, your ability to apply methods, and your engagement with the material.
\begin{description}
  \item[Weekly Problem Sets (20\%):] Short assignments focusing on theoretical concepts, DAG interpretation, and articulating assumptions.
  \item[Applied R Assignments (25\%):] A series of five assignments where you will independently apply methods to real or synthetic datasets.
  \item[Midterm Exam (20\%):] Covers foundational concepts from the first half of the course (Weeks 1-6).
  \item[Final Project (35\%):] A capstone project applying causal methods to a dataset. The project includes:
    \begin{itemize}
      \item[] \textbf{Project Proposal (5\% of total grade):} A clear causal question, dataset identification, and an initial DAG. Due mid-semester.
      \item[] \textbf{Final Report \& Presentation (30\% of total grade):} A comprehensive report detailing your methods, findings, and limitations, accompanied by a short in-class presentation. Due during finals week.
    \end{itemize}
\end{description}

\vspace{6mm}
\subsection*{Grading Scale}
Final letter grades will be assigned based on the following scale:
\begin{center}
  \begin{tabular}{ll@{\hskip 2em}ll}
    A+ & 97-100\% & B- & 80-82.9\% \\
    A  & 93-96.9\% & C+ & 77-79.9\% \\
    A- & 90-92.9\% & C  & 73-76.9\% \\
    B+ & 87-89.9\% & C- & 70-72.9\% \\
    B  & 83-86.9\% & F  & < 70\% \\
  \end{tabular}
\end{center}

\vspace{6mm}
\section*{Recommended Materials}
\begin{description}
  \item[Primary Texts:]
    \begin{itemize}
      \item Hernán MA, Robins JM. \textit{Causal Inference: What If.} (\href{https://www.hsph.harvard.edu/miguel-hernan/causal-inference-book/}{Online})
      \item Pearl J, et al. \textit{Causal Inference in Statistics: A Primer.}
      \item Imbens GW, Rubin DB. \textit{Causal Inference for Statistics, Social, and Biomedical Sciences.}
    \end{itemize}
  \item[Supplementary Texts:]
    \begin{itemize}
      \item VanderWeele TJ. \textit{Explanation in Causal Inference.}
      \item Brumback BA. \textit{Fundamentals of Causal Inference with R.}
    \end{itemize}
  \item[Software:] We will use \textbf{R} and \textbf{RStudio}. Key packages will include \texttt{dagitty}, \texttt{MatchIt}, \texttt{WeightIt}, \texttt{cobalt}, \texttt{ivreg}, and \texttt{mediation}.
\end{description}

\vspace{6mm}
\section*{Course Schedule}
\begin{center}
  \small
  \begin{longtable}{p{0.1\textwidth}p{0.2\textwidth}p{0.6\textwidth}}
    \toprule
    \textbf{Week} & \textbf{Topic} & \textbf{Key Concepts and Readings} \\
    \midrule
    \endhead
    \endhead
    1 & Introduction & Association vs. Causation, Simpson's Paradox, Intro to Potential Outcomes \& \textit{Readings:} HR Ch 1; PGJ Ch 1 \\
    \midrule
    2 & Potential Outcomes & ATE/ATT, SUTVA, Consistency, Randomized Experiments \& \textit{Readings:} HR Ch 2-3; IR Ch 2-3 \\
    \midrule
    3 & Causal Graphs (DAGs) & Nodes, Edges, Paths, d-separation, Backdoor Criterion \& \textit{Readings:} PGJ Ch 2; HR Ch 6-7 \\
    \midrule
    4 & Bias and Identification & Colliders, M-Bias, Selection Bias, Frontdoor Criterion \& \textit{Readings:} PGJ Ch 3, HR Ch 8 \\
    \midrule
    5 & G-Formula & Conditional Exchangeability, Positivity, Standardization \& \textit{Readings:} HR Ch 4-5 \\
    \midrule
    6 & Propensity Scores & Balancing Scores, Matching, Stratification, IPTW \& \textit{Readings:} HR Ch 12-13; IR Ch 13-18 \\
    \midrule
    7 & \textbf{Midterm Exam} & Covers material from Weeks 1-6. \\
    \midrule
    8 & Instrumental Variables & Exclusion Restriction, LATE, 2SLS\@. \textbf{Project Proposal Due\@.} \& \textit{Readings:} HR Ch 16; IR Ch 23 \\
    \midrule
    9 & Causal Mediation I & NDE, NIE, Counterfactual Definitions, Identification Assumptions \& \textit{Readings:} VanderWeele Ch 1--2 \\
    \midrule
    10 & Causal Mediation II & Estimation (Difference/Product Methods), Interaction \& \textit{Readings:} VanderWeele Ch 3--4 \\
    \midrule
    11 & Sensitivity Analysis & E-values, Bounding Factors, Robustness to Unmeasured Confounding \& \textit{Readings:} HR Ch 14; Ding \& VanderWeele (2016) \\
    \midrule
    12 & Marginal Structural Models & Time-Varying Confounding, IPW for MSMs, Sequential Exchangeability \& \textit{Readings:} HR Ch 17-18, 21 \\
    \midrule
    13 & Project Work & In-class workshop, final project support, and Q\&A. \\
    \midrule
    \bottomrule
  \end{longtable}
  \normalsize
\end{center}

\vspace{6mm}
\section*{Course Policies}
\begin{description}
  \item[Late Work Policy:] Deadlines are firm. Late assignments will be penalized 10\% for each day they are late. Exceptions will only be granted for documented emergencies, and you must contact the instructor as soon as possible, preferably before the deadline.
  \item[Academic Integrity:] All work submitted in this course must be your own. Collaboration on high-level concepts for problem sets is encouraged, but all submitted code and written answers must be your own original work. Copying code from classmates or online sources without attribution is a serious violation of academic integrity. Please consult the university's academic integrity policy for detailed information.
  \item[Communication:] The primary method for course announcements will be via the university's learning management system. For questions about course content, please use office hours or ask in class. For personal or administrative matters, email the instructor; please allow 24-48 business hours for a response.
  \item[Classroom Environment:] This course is designed to be an inclusive and welcoming environment for all students. Diverse viewpoints are valued and encouraged. All students are expected to contribute to a respectful and collaborative atmosphere during class discussions.
  \item[Accessibility:] Students with disabilities who may need accommodations in this class are encouraged to contact the university's Accessibility Resource Center as soon as possible to ensure that such accommodations are arranged in a timely fashion.
\end{description}

\end{document}
